\subsection{极大交换子代数}

我们称 K-代数 A 的一个子代数是 A 的极大交换子代数，如果它是 A 的交换子代数所成之集中的一个极大元素。

引理 3. —— 设 A 是一个 K-代数，L 是 A 的一个子代数。

\begin{enumerate}
\def\labelenumi{\alph{enumi})}
\item
  代数 \(\mathrm{L}\) 是 \(\mathrm{A}\) 的极大交换子代数当且仅当 \(\mathrm{L}\) 等于它的交换子 \(\mathrm{L}'\)（在 \(\mathrm{A}\) 中的）。
\item
  设 \(\mathrm{K}'\) 是一个非零的交换 K-代数。则 \(\mathrm{L}\) 是 A 的极大交换子代数当且仅当 \(\mathrm{L}_{(\mathrm{K}')}\) 是 \(\mathrm{A}_{(\mathrm{K}')}\) 的极大交换子代数。
\end{enumerate}

我们来证明 a)。首先假设 L 等于 \(L'\)。那么 L 是交换的；若 M 是 A 的一个包含 L 的交换子代数，则我们有 \(xy = yx\)（对 \(x\) 属于 \(\mathrm{L}\) 与 \(y\) 属于 \(\mathrm{M}\)），于是 \(\mathrm{M} \subset \mathrm{L}'\)，从而 \(\mathrm{M} = \mathrm{L}\)。因此，\(\mathrm{L}\) 是 \(\mathrm{A}\) 的极大交换子代数。

反之，假设 L 是 A 的一个极大交换子代数，并设 x 是 \(L'\) 的一个元素。则由 \(L \cup \{x\}\) 生成的 A 的子代数 M 是交换的且包含 L。由于 L 是极大的，我们有 M = L，于是 \(x \in L\)，最终 \(L = L'\)，这就给出 a)。

由 III, §4, No.~4, 第 468 页的命题 6，\(\mathrm{L}_{(\mathrm{K}^{\prime})}\) 在 \(\mathrm{A}_{(\mathrm{K}^{\prime})}\) 中的交换子是 \(\mathrm{L}'_{(\mathrm{K}^{\prime})}\)。由于等式 \(\mathrm{L} = \mathrm{L}'\) 与 \(\mathrm{L}_{(\mathrm{K}^{\prime})} = \mathrm{L}'_{(\mathrm{K}^{\prime})}\) 等价（II, §7, No.~9, 第 311 页，命题 19），断言 b) 由 a) 得出。

命题 3. —— 设 A 是一个有限次数的中心单 K-代数，L 是 A 的一个半单交换子代数。下列性质等价：

\begin{enumerate}
\def\labelenumi{(\roman{enumi})}
\item
  代数 \(\mathrm{L}\) 是 \(\mathrm{A}\) 的一个极大交换子代数。
\item
  左 L-模 A 是自由的，维数为 {[}L : K{]}。
\item
  我们有 \([A:K]=[L:K]^{2}\)。
\end{enumerate}

此外，假设 A 是代数 \(\operatorname{End}_{\mathrm{K}}(\mathrm{V})\)，其中 V 是 K 上一个非零有限维的向量空间。则前面的性质也与下面的性质等价：

\begin{enumerate}
\def\labelenumi{(\roman{enumi})}
\setcounter{enumi}{3}
\tightlist
\item
  V 是一个维数为 1 的自由 L-模。
\end{enumerate}

\begin{enumerate}
\def\labelenumi{\Alph{enumi})}
\tightlist
\item
  假设 A 具有形式 \(\operatorname{End}_{\mathrm{K}}(\mathrm{V})\)，其中 V 是域 K 上一个非零有限维的向量空间。我们将按照下图来建立条件 (i) 至 (iv) 的等价性：
\end{enumerate}

\[
(i) \Longrightarrow (i v) \Longrightarrow (i i) \Longrightarrow (i i i) \Longrightarrow (i).
\]

由于 L 是 K 上一个有限次数的半单交换代数，我们可以把它等同为 K 的有限次数扩张的一个有限积 \(\prod_{i\in I}L_{i}\)（VIII, 第 137 页，命题 3）。对每个 \(i\in I\)，用 \(V_{i}\) 表示 L-模 V 的 \(L_{i}\) 型同型分量；由于 V 是一个忠实 L-模，它是 \(L_{i}\) 上一个非零有限维的向量空间（VIII, 第 144 页，推论）。于是我们可以把 V 等同为 \(\prod_{i\in I}V_{i}\)。在这些条件下，L 在 A 中的交换子 \(L'\)（它就是代数 \(\operatorname{End}_{\mathrm{L}}(\mathrm{V})\)）可以等同为积 \(\prod_{i\in I}\operatorname{End}_{\mathrm{L}_{i}}(\mathrm{V}_{i})\)。

假设 L 是 \(\operatorname{End}_{\mathrm{K}}(\mathrm{V})\) 的一个极大交换子代数。由引理 3, a)，我们有 \(L = L'\)，于是 \(L'\) 是交换的，且有 \(\dim_{L_i}(V_i) = 1\)（对每个 \(i \in I\)）。因此 (i) 蕴含 (iv)。

假设 L-模 V 是维数为 1 的自由模。设 \((e_{1},\ldots,e_{r})\) 是 V 在 K 上的一个基。映射 \(a\mapsto(ae_{1},\ldots,ae_{r})\) 是一个左 A-模同构，因而是从 A 到 \(V^{r}\) 的 L-模同构。于是，A 是一个维数为 \(r\) 的自由左 L-模，且我们有 \(r = \dim_{\mathrm{K}}(\mathrm{V}) = [\mathrm{L}:\mathrm{K}]\)。因此 (iv) 蕴含 (ii)。

显然 (i) 蕴含 (iii)。

最后，假设我们有 \([A:K]=[L:K]^{2}\)，换言之，\(\dim_{\mathrm{K}}(\mathrm{V})=[\mathrm{L}:K]\)。于是我们有 \(\sum_{i}\dim_{\mathrm{L}_{i}}(\mathrm{V}_{i})[\mathrm{L}_{i}:\mathrm{K}]=\sum_{i}[\mathrm{L}_{i}:\mathrm{K}]\)，从而对每个 i，\(V_{i}\) 在 \(L_{i}\) 上的维数为 1。于是对每个 i 有 \(\operatorname{End}_{\mathrm{L}_{i}}(\mathrm{V}_{i})=\mathrm{L}_{i}\)，故 \(L^{\prime}=L\)。由引理 3, a)，L 是 A 的一个极大交换子代数。因此 (iii) 蕴含 (i)。

\begin{enumerate}
\def\labelenumi{\Alph{enumi})}
\setcounter{enumi}{1}
\tightlist
\item
  我们继续讨论一般情形。由 VIII, 第 252 页的定理 1，存在一个可分扩张 \(\mathrm{K}'\)（\(\mathrm{K}\) 的、有限次数的），使得 \(\mathrm{K}'\)-代数 \(\mathrm{A}_{(\mathrm{K}')}\) 同构于一个代数 \(\mathrm{End}_{\mathrm{K}'}(\mathrm{V}')\)，其中 \(\mathrm{V}'\) 是 \(\mathrm{K}'\) 上一个非零有限维的向量空间。于是 \(\mathrm{K}'\)-代数 \(\mathrm{L}_{(\mathrm{K}')}\) 是交换的且是半单的（VIII, 第 222 页，推论 2）。由证明的第一部分，下列性质等价：
\end{enumerate}

(i') 代数 \(L_{(K')}\) 是 \(A_{(K')}\) 的一个极大交换子代数。

(ii') 左 \(\mathrm{L}_{(\mathrm{K}^{\prime})}\)-模 \(\mathrm{A}_{(\mathrm{K}^{\prime})}\) 是自由的，维数为 \([\mathrm{L}_{(\mathrm{K}^{\prime})}:\mathrm{K}^{\prime}]\)。

(iii') 我们有 \(\left[\mathrm{A}_{\left(\mathrm{K}^{\prime}\right)}:\mathrm{K}^{\prime}\right] = \left[\mathrm{L}_{\left(\mathrm{K}^{\prime}\right)}:\mathrm{K}^{\prime}\right]^{2}\)。

由引理 3, b)，性质 (i) 与 (i') 等价。

设 \(n = [\mathrm{L}:\mathrm{K}]\)；则 \(n = [\mathrm{L}_{(\mathrm{K}^{\prime})}:\mathrm{K}^{\prime}]\)。性质 (ii) 表示左 L-模 A 与 \(\mathrm{L}^n\) 同构；由 VIII, 第 37 页的定理 3，这等价于 \(\mathrm{L}_{(\mathrm{K}^{\prime})}\)-模 \(\mathrm{A}_{(\mathrm{K}^{\prime})}\) 与 \((\mathrm{L}_{(\mathrm{K}^{\prime})})^n\) 同构。由此得到 (ii) 与 (ii') 的等价性。

最后，我们有 \([A:K]=[A_{(K')}:K']\) 与 \([L:K]=[L_{(K')}:K']\)，由此给出性质 (iii) 与 (iii') 的等价性。我们这样就证明了性质 (i)、(ii) 与 (iii) 的等价性。

推论. —— 设 A 是 K 上一个有限次数的中心单代数，L 是一个半单交换 K-代数，使得 \([A:K]\) 等于 \([L:K]^{2}\)。设 f 与 g 是从 L 到 A 的单同态。存在 A 的一个内自同构 \(\theta\)，使得 \(g = \theta \circ f\)。

设 \(n = [L : K]\)。作为子环 \(f(L)\) 上的左模，A 是维数为 n 的自由模；这由命题 3 的性质 (ii) 与 (iii) 的等价性得出。由于 f 是从 L 到 \(f(L)\) 的同构，左 L-模 \(A^{f}\)（其作用律由 \((x, a) \mapsto f(x)a\) 给出）是维数为 n 的自由模。对 \(A^{g}\) 同样如此，因而它同构于 \(A^{f}\)。我们利用命题 1 的性质 (i) 与 (ii) 的等价性（VIII, 第 257 页）得出结论。

假设 \(\mathbf{A}\) 是 \(\mathbf{K}\) 上一个有限次数的中心单代数。

可以存在 A 的极大交换子代数 L，它不是半单的，且使得 \([A:K]\neq[L:K]^{2}\)（VIII, 第 270 页，习题 5）。

